\documentclass[10pt,a4paper]{amsart}
\usepackage[margin=19mm]{geometry}
\usepackage{amsmath,amssymb,mathtools}
\usepackage[scr=rsfs,cal=euler]{mathalfa}
\usepackage{xfrac}
\usepackage[unicode,hidelinks,bookmarks=false]{hyperref}
\newcommand{\ie}{i.e.\ }
\newcommand{\cf}{cf.\ }
\newcommand{\resp}{respectively\ }
\newcommand{\Subsection}{\subsection}
\providecommand{\nobreakdash}{\nobreak\hbox{-}}
\setlength{\emergencystretch}{2em}
\multlinegap0pt
\usepackage{bbm}
\usepackage{amssymb}
\usepackage{mathtools}
\usepackage[all]{xy}
\usepackage{tikz}
\usepackage{enumerate}
\usepackage[shortlabels]{enumitem}
\newtheorem{theorem}{Theorem}
\title{Extremal Graphs for the Lights Out Problem}
\author{Julien Codsi, Sergio Cristancho, Alexander Divoux, Varun Sivashankar}
\date{Source: arXiv:2602.07241v1 (CC BY 4.0)}
\begin{document}
\maketitle

\noindent\emph{Source and licence.} Extremal Graphs for the Lights Out Problem. Version arXiv:2602.07241v1 (CC BY 4.0), distributed by arXiv under the Creative Commons Attribution 4.0 International licence, http://creativecommons.org/licenses/by/4.0/, as recorded in the arXiv licence metadata for this exact version and shown on the abstract page https://arxiv.org/abs/2602.07241v1. Full licence notice: \texttt{licenses/arxiv-2602.07241-CC-BY-4.0.txt}.\par\smallskip

\noindent\textit{Editorial scope.} Counted unit: the article's theorem thm:oddset (source0.tex lines 310--312). The article's complete original proof of it is delivered with it (source0.tex lines 313--349, one \texttt{proof} environment of 37 lines). No mathematics is written, repaired or re-derived: the statement and the proof are the article's own lines, in the article's own notation, and no retained line is substituted or re-flowed.  No further source-local context is needed: every cross-reference of every retained line resolves inside the two delivered spans.  Nothing was omitted from any retained span, so the package carries no omission record.  No retained line contains a citation, so the wrapper carries no bibliography and no bibliography entry.  No retained line contains a figure environment, an image-inclusion command or drawing code, so no image is delivered and the package declares no asset dependency.  Editorial formatting only, in addition: an AMS \texttt{amsart} preamble that restates the article's own declarations and macros used by the retained lines, namely the environments \texttt{theorem} on separate counters, together with the standard packages those lines need.  The retained environments print in this document as Theorem 1 in order of appearance, whereas the article prints the same items with the numbering of its own full document.  The complete native source of the article as distributed in the arXiv e-print for this exact version is bundled unchanged as \texttt{sources/194164/source0.tex}.
\begin{theorem}\label{thm:oddset}
    Let $G$ be an even graph and $S\subseteq V(G)$ with $|S|$ odd. Then for any $k$, the number of size-$k$ matchings covering $S$ is even. 
\end{theorem}

\begin{proof}
    Let $G = (V, E)$, and work in the ring $R = \mathbb F_2[t][x_v:v\in V]/\{x_v^2:v\in V\}$. Define the polynomial $P = \prod_{ab\in E}(1+tx_ax_b)\in R$. For $A\subseteq V$, write $x^A = \prod_{a\in A}x_a$. By the definition of $R$, we have
    \[
        P = \sum_{M\text{ matching}}t^{|M|}x^{V(M)} = \sum_{i\ge 0}\sum_{\substack{U\subseteq V\\|U|=2i}} c_Ut^ix^U,
    \]
    where $c_U$ is the number of perfect matchings of $G[U]$ modulo $2$. Let $T = V\setminus S$ and define a second polynomial $Q = \prod_{w\in T}(1+x_w)\in R$. Consider the coefficient $[t^kx^V](P Q)$. Any $U\subseteq V$ with $|U| \neq 2k$ cannot contribute to this coefficient as its monomial misses $t^k$. Any $U\subseteq V$ with $U\not\supseteq S$ cannot contribute to this coefficient as its monomial misses $x_s$ for some $s\in S$. Any $U\subseteq V$ with $|U|=2k$ and $U\supseteq S$ contributes $c_U$ to this coefficient exactly once: choose ``$1$'' for $w\in T\cap U$ and ``$x_w$'' for $w\in T\setminus U$ when multiplying by $Q$. Therefore,
    \[
        [t^kx^V](P Q) = \sum_{\substack{U\subseteq V\\|U|=2k\\U\supseteq S}} c_U = m_k(G, S),
    \]
    where $m_k(G, S)$ is the number of size-$k$ matchings covering $S$ modulo $2$. For each $v\in V$, define the differential operator $\partial_v$ with $\partial_v(x_v) = 1$, $\partial_v(x_u) = 0$ for $u\neq v$, and $\partial_v(t) = 0$. Define $\partial = \sum_{v\in V}\partial_v$. For each $v\in V$, using the fact that $(1+tx_u x_v)^2 = 1$, the product rule gives
    \[
        \partial_v P = \partial_v\prod_{ab\in E}(1+tx_ax_b) = \sum_{u\in N(v)}tx_uP(1+tx_ux_v) = tP\sum_{u\in N(v)}x_u.
    \]
    Summing over all $v\in V$ and using $\deg(v) = 0$ modulo $2$ as $G$ is even,
    \[
        \partial P = tP\sum_{v\in V}\sum_{u\in N(v)}x_u = tP\sum_{v\in V}\deg(v)x_v = 0.
    \]
    Now consider the coefficient $[t^kx^V]\partial(P Q)$. Since $\partial$ reduces the degrees of all monomials by $1$, it must be that $[t^kx^V]\partial(P Q) = 0$. On the other hand, expanding $\partial(P Q)$ with the product rule and using $\partial P=0$,
    \[
        \partial(P Q) = (\partial P) Q + P(\partial Q) = P(\partial Q).
    \]
    Since $Q = \prod_{w\in T}(1+x_w)$, we have $\partial Q  = Q\sum_{w\in T}(1+x_w) = Q(|T| + \sum_{w\in T}x_w)$. Hence $\partial  (P Q) = P Q(|T| + \sum_{w\in T}x_w)$. Gathering coefficients, this gives
    \[
        [t^kx^V]\partial(P Q) = |T|m_k(G, S) + \sum_{w\in T}[t^kx^{V\setminus\{w\}}](P Q).
    \]
    Now fix $w\in T$ and consider the coefficient $[t^kx^{V\setminus\{w\}}](PQ)$. The same reasoning for determining $[t^k x^V](PQ)$ yields that any $U\subseteq V$ with $|U| \neq 2k$ or $U\not\supseteq S$ cannot contribute to this coefficient. Furthermore, any $U$ with $w\in U$ also cannot contribute, since it cannot have $x_w$ in its monomial. Finally, any $U\subseteq V$ with $|U|=2k$, $U\supseteq S$, and $U\not\ni w$ contributes $c_U$ to this coefficient exactly once: choose ``$1$'' for the vertices in $T\cap U$ and $w$, and choose $x_a$ for $a\in T\setminus U$ when multiplying by $Q$. Therefore, 
    \[
        [t^kx^{V\setminus\{w\}}](P Q) = \sum_{\substack{U\subseteq V\\|U|=2k\\U\supseteq S\\U\not\ni w}} c_U
    \]
    is the number of size-$k$ matchings covering $S$ but not covering $w$. Summing over all $w\in T$ and interpreting the scalars $|S|$, $|T|$, $|U|$, and $|V|$ modulo $2$,
        \begin{align*}
        \sum_{w\in T}[t^kx^{V\setminus\{w\}}](P Q) &= \sum_{w\in T}\sum_{\substack{U\subseteq V\\|U|=2k\\U\supseteq S\\U\not\ni w}} c_U = \sum_{w\in T}\sum_{\substack{U\subseteq V\\|U|=2k\\U\supseteq S}} \mathbbm{1}[w\notin U]c_U\\
        &= \sum_{\substack{U\subseteq V\\|U|=2k\\U\supseteq S}}c_U|T\setminus U| = \sum_{\substack{U\subseteq V\\|U|=2k\\U\supseteq S}}c_U(|V|-|U|)\\
        &= \sum_{\substack{U\subseteq V\\|U|=2k\\U\supseteq S}}c_U(|V|-2k) = |V|\sum_{\substack{U\subseteq V\\|U|=2k\\U\supseteq S}}c_U = |V|m_k(G, S).
    \end{align*}
    It follows that $0 = [t^kx^V]\partial(P Q) = |T|m_k(G, S) + |V|m_k(G, S) = (|V|-|T|)m_k(G,S) = |S|m_k(G,S)$. Since $|S|$ is odd, this means $m_k(G,S) = 0$, so the number of size-$k$ matchings covering $S$ is even.
\end{proof}

\end{document}
